\section{Introduction}

Durga is a BPMN-driven Kafka scaffolding tool. It reads BPMN models, extracts
an executable subset of process semantics, and generates Java and Kafka-oriented
project skeletons that run as production-grade applications on Kafka.

The project is a generator, not a workflow engine:
\begin{itemize}
\item BPMN models are parsed with Camunda (but not run in Camunda),
\item worker, gateway, orchestrator and event handler classes are rendered
      with StringTemplate from the process graph,
\item Kafka topics and reactive-messaging bindings are generated alongside
      the code,
\item helper scripts and probes are generated so the resulting process can be
      exercised and observed,
\item a separate Kafka Streams monitoring topology materializes process state
      and operational views.
\end{itemize}

This manual describes the current system. It focuses on what Durga generates,
what runtime model those artifacts implement, and where the support boundary ends.

The document is organized in that order. Section~\ref{sec:SystemOverview} describes the major subsystems. Section~\ref{sec:ExecutionModel} explains the BPMN-to-Kafka execution mapping. Section~\ref{sec:Monitoring} describes the canonical process-event and monitoring model. Section~\ref{sec:GeneratedProjects} describes the generated project output, and Section~\ref{sec:CurrentStatus} summarizes maturity, constraints and current gaps.
